
\documentclass{article}
\usepackage{colortbl}
\usepackage{makecell}
\usepackage{multirow}
\usepackage{supertabular}

\begin{document}

\newcounter{utterance}

\centering \large Interaction Transcript for game `cladder', experiment `full\_v1.5\_default', episode 9813 with refineonly{-}Qwen3.5{-}4B.
\vspace{24pt}

{ \footnotesize  \setcounter{utterance}{1}
\setlength{\tabcolsep}{0pt}
\begin{supertabular}{c@{$\;$}|p{.15\linewidth}@{}p{.15\linewidth}p{.15\linewidth}p{.15\linewidth}p{.15\linewidth}p{.15\linewidth}}
    \# & \multicolumn{2}{c}{Player} && \multicolumn{2}{c}{Game Master} \\
    \hline

    \theutterance \stepcounter{utterance}  
    & & & \multicolumn{4}{p{0.6\linewidth}}{
        \cellcolor[rgb]{0.9,0.9,0.9}{
            \makecell[{{p{\linewidth}}}]{
                \texttt{\tiny{[P1$\langle$GM]}}
                \texttt{You are an expert in causal inference. The following question is not a typical commonsense query, but rather a meticulously designed question created by a professor specializing in causal inference, intended to assess the students' mastery of the course content.} \\
\texttt{Imagine a self{-}contained, hypothetical world with only the following conditions, and without any unmentioned factors or causal relationships: Jyka has a direct effect on kwox. Yupt has a direct effect on kwox. We know that jyka or yupt causes kwox. We observed an individual is not yupt. Would an individual is kwox if jyka instead of not jyka?} \\
\texttt{Start your answer with 'Yes' or 'No', followed by additional reasoning or evidence to support your explanation.} \\
            }
        }
    }
     \\ \\

    \theutterance \stepcounter{utterance}  
    & & \multicolumn{4}{p{0.6\linewidth}}{
        \cellcolor[rgb]{0.95,0.95,0.95}{
            \makecell[{{p{\linewidth}}}]{
                \texttt{\tiny{[P1$\leftrightarrow$P]}}
                \texttt{{[}DRAFT{]}} \\
\texttt{Yes} \\
\\ 
\texttt{**Reasoning:**} \\
\\ 
\texttt{Based on the conditions provided in this hypothetical world, we can deduce the outcome using the principles of causal sufficiency and the absence of confounding variables.} \\
\\ 
\texttt{1.  **Causal Structure**: We are given two direct causal pathways: \$Jyka \textbackslash{}rightarrow Kwox\$ and \$Yupt \textbackslash{}rightarrow Kwox\$.} \\
\texttt{2.  **Sufficiency Condition**: The prompt states, "We know that jyka or yupt causes kwox." In the context of causal inference problems of this type, this phrasing implies a **sufficient cause** relationship. It means that for an individual to be \$Kwox\$, at least one of the causes (\$Jyka\$ or \$Yupt\$) must be present. Conversely, if neither cause is present, the effect cannot occur.} \\
\texttt{3.  **Current State**: We observe an individual who is **not** \$Yupt\$.} \\
\texttt{4.  **Hypothetical Intervention**: We are asked if the individual would be \$Kwox\$ if they were \$Jyka\$ (instead of not \$Jyka\$).} \\
\texttt{5.  **Logical Deduction**:} \\
\texttt{    *   If the individual is \$Jyka\$, the condition "Jyka or Yupt" is satisfied (since \$Jyka\$ is true).} \\
\texttt{    *   Since the world contains *only* the mentioned factors and causal relationships, there are no other unmentioned causes required for \$Kwox\$.} \\
\texttt{    *   Therefore, the presence of \$Jyka\$ is sufficient to trigger the effect \$Kwox\$.} \\
\\ 
\texttt{Thus, if the individual is \$Jyka\$, they will be \$Kwox\$.} \\
            }
        }
    }
    & \\ \\

    \theutterance \stepcounter{utterance}  
    & \multicolumn{4}{p{0.6\linewidth}}{
        \cellcolor[rgb]{0.561,0.737,0.561}{
            \makecell[{{p{\linewidth}}}]{
                \texttt{\tiny{[P1$\rangle$GM]}}
                \texttt{Yes} \\
\\ 
\texttt{**Reasoning:**} \\
\\ 
\texttt{Based on the conditions provided in this hypothetical world, we can deduce the outcome using the principles of causal sufficiency and the absence of confounding variables.} \\
\\ 
\texttt{1.  **Causal Structure**: We are given two direct causal pathways: \$Jyka \textbackslash{}rightarrow Kwox\$ and \$Yupt \textbackslash{}rightarrow Kwox\$.} \\
\texttt{2.  **Sufficiency Condition**: The prompt states, "We know that jyka or yupt causes kwox." In the context of causal inference problems of this type, this phrasing implies a **sufficient cause** relationship. It means that for an individual to be \$Kwox\$, at least one of the causes (\$Jyka\$ or \$Yupt\$) must be present. Conversely, if neither cause is present, the effect cannot occur.} \\
\texttt{3.  **Current State**: We observe an individual who is **not** \$Yupt\$.} \\
\texttt{4.  **Hypothetical Intervention**: We are asked if the individual would be \$Kwox\$ if they were \$Jyka\$ (instead of not \$Jyka\$).} \\
\texttt{5.  **Logical Deduction**:} \\
\texttt{    *   If the individual is \$Jyka\$, the condition "Jyka or Yupt" is satisfied (since \$Jyka\$ is true).} \\
\texttt{    *   Since the world contains *only* the mentioned factors and causal relationships, there are no other unmentioned causes required for \$Kwox\$.} \\
\texttt{    *   Therefore, the presence of \$Jyka\$ is sufficient to trigger the effect \$Kwox\$.} \\
\\ 
\texttt{Thus, if the individual is \$Jyka\$, they will be \$Kwox\$.} \\
            }
        }
    }
    & & \\ \\

    \theutterance \stepcounter{utterance}  
    & & & \multicolumn{2}{p{0.3\linewidth}}{
        \cellcolor[rgb]{0.9,0.9,0.9}{
            \makecell[{{p{\linewidth}}}]{
                \texttt{\tiny{[GM$|$GM]}}
                \texttt{Yes} \\
            }
        }
    }
    & & \\ \\

    \theutterance \stepcounter{utterance}  
    & & & \multicolumn{2}{p{0.3\linewidth}}{
        \cellcolor[rgb]{0.9,0.9,0.9}{
            \makecell[{{p{\linewidth}}}]{
                \texttt{\tiny{[GM$|$GM]}}
                \texttt{yes} \\
            }
        }
    }
    & & \\ \\

    \theutterance \stepcounter{utterance}  
    & & & \multicolumn{2}{p{0.3\linewidth}}{
        \cellcolor[rgb]{0.9,0.9,0.9}{
            \makecell[{{p{\linewidth}}}]{
                \texttt{\tiny{[GM$|$GM]}}
                \texttt{game\_result = WIN} \\
            }
        }
    }
    & & \\ \\

\end{supertabular}
}

\end{document}
